\documentclass[10pt,a4paper]{article}
\usepackage[a4paper,margin=2.2cm]{geometry}
\usepackage[T1]{fontenc}
\usepackage[utf8]{inputenc}
\usepackage{amsmath,amssymb}
\usepackage{booktabs,array,graphicx}
\usepackage{microtype}
\pagestyle{empty}
\setlength{\parindent}{0pt}
\renewcommand{\arraystretch}{1.12}
\newcommand{\celula}[1]{\begin{tabular}[c]{@{}c@{}}#1\end{tabular}}
\newsavebox{\tabcaixa}
\newcommand{\tabajusta}[1]{\sbox{\tabcaixa}{#1}\ifdim\wd\tabcaixa>\linewidth\resizebox{\linewidth}{!}{\usebox{\tabcaixa}}\else\usebox{\tabcaixa}\fi}
\begin{document}

{\large\bfseries Table S3. Trained models at $g = 5$~km in the four Bauru cells (five draws)}

\medskip
\small
A trained model's valid-node MAE drifts with the split draw. Split draws $101$--$105$ in the four Bauru cells, $g = 5$~km, $b = 2$~km, training seed $42$, training budget of eight epochs. Valid-node mean absolute error of received power, in decibels, mean over draws (between-draw standard deviation) for the graph regressor (GNN) and the perceptron (MLP); bottom row: standard deviation of the per-cell means. Over the five draws, the between-draw standard deviation of the valid-node MAE is $5.2$ to $21.6$ times (GNN) and $4.7$ to $32.0$ times (MLP) the same-seed repeat difference ($0.132$~dB, the largest difference between two training repeats with the same seed and partition, measured at $5$~km). In Bauru Q1 and Bauru Q3, where five training seeds were also run at split seed $42$, it is $11.7$ (GNN) and $38.2$ (MLP) times the standard deviation across the five training seeds in Bauru Q1, and $2.8$ (GNN) and $50.6$ (MLP) times in Bauru Q3.

\medskip
\centering
\tabajusta{\begin{tabular}{lcc}
\toprule
Cell & \celula{GNN valid \\ \relax MAE $\downarrow$ (dB): \\ \relax mean (SD \\ \relax over draws)} & \celula{MLP valid \\ \relax MAE $\downarrow$ (dB): \\ \relax mean (SD \\ \relax over draws)} \\
\midrule
Bauru Q1 & 1.80 (1.25) & 3.81 (1.71) \\
Bauru Q2 & 2.90 (2.85) & 6.03 (4.22) \\
Bauru Q3 & 1.61 (0.98) & 3.09 (2.93) \\
Bauru Q4 & 1.11 (0.69) & 1.57 (0.62) \\
\midrule
SD across the four cells of the per-cell means & 0.76 & 1.85 \\
\bottomrule
\end{tabular}
}

\end{document}
